I dette prosjektet har vi funnet opp, utviklet og evaluert et alternativt spillkonsept. Et alternativt spill defineres som \textit{“et spill hvor en eller flere av karakteristikkene til spillet avviker fra et vanlig spill.”} Prosjektet bygger på en forstudie der eksisterende alternative spill og grunnleggende spilldesignteori ble studert.

Konseptet \textit{Hover} er et gripende spill hvor lokasjonen brukes som den primære måte å kontrollere spillet på. Spillere får poeng for å være på bestemte steder i den virkelige verden som kan forbedre deres mentalitet. Eksempler på steder inkluderer treningssentre, skoler og museer. Sosiale elementer og andre typiske spilldesignelementer har blitt brukt til å utvikle et engasjerende spill.

Hover ble testet på to grupper mennesker for å evaluere hvordan spillet påvirket spillernes engasjement, motivasjon, glede og vaner. Forskjeller mellom gruppene ble også vurdert. Deltakerne brukte applikasjonen i to uker før de ble bedt om å fylle ut et spørreskjema. Svarene fra dette spørreskjemaet, i tillegg til observasjoner av deltakerne underveis i testingen, intervjuer og analysedata, ble brukt til å konkludere studien.

Testingen avslørte at spillets sosiale elementer påvirket spillernes engasjement, motivasjon og glede. Andre spillelementer, som poeng og toppliste, påvirket også spillernes oppfatning av spillet. Derimot ble ikke spillernes vaner påvirket i nevneverdig grad. Eksperimentet illustrerte også at det sosiale båndet mellom deltakerne i en gruppe påvirket bruken av applikasjonen.

Selv om den pågående Covid-19-pandemien hadde stor innvirkning på testingen av Hover, viser resultatene lovende potensial. Dermed anbefales videreutvikling og testing av konseptet for å kunne trekke en fullstendig konklusjon.

\par\vspace*{\fill} % Moves keywords to the bottom of the page
\textbf{\textit{Nøkkelord --}} NTNU, Datateknologi, Programvareutvikling, Spillutvikling, Masterprosjekt, Alternative spill, Gripende spill, Gamification, Samfunnsnyttig, Lokasjonsbasert % Add all the keywords associated with your thesis here